% Esquisse d'un programme (cote 119, copy 1): archive cote 119, site page 7 (PDF page 8), typescript page 5.
% typescript. Draft compared: EsquisseFr.pdf (text layer):170-231.
% AI-transcribed by Claude Opus 5.5 (claude-opus-5-5, adaptive thinking, effort high), Message Batches API, batch esq1 (msgbatch_01LNsogFyF3y6KSp9UbtJtBk), 2026-11-12; not yet checked by eye.
\scanpage{119/7}
\mspage{5}

de $\hat{\pi}_{0,3}$ i.e. sur le ou les paramètres correspondants, pour qu'il
provienne d'un élément de $I\!\Gamma$ - ce qui fournirait une description
``purement algébrique'', en termes de groupes profinis et sans ré-%
férence à la théorie de Galois des corps de nombres, du groupe de
Galois $I\!\Gamma = \mathrm{Gal}(\bar{\underline{Q}}/\underline{Q})$ !

Peut-être une caractérisation même conjecturale de $I\!\Gamma$
comme sous-groupe de Autext($\hat{\pi}_{0,3}$) est-elle pour le moment hors de
portée($^1$) ; je n'ai pas de conjecture à proposer encore. Une autre
tâche par contre est abordable immédiatement, c'est celle de décrire
l'action de $I\!\Gamma$ sur toute la tour de Teichmüller, en termes de son
action sur le ``premier étage'' $\hat{\pi}_{0,3}$ , i.e. exprimer un automorphisme
de cette tour, en termes du ``paramètre'' dans $\hat{\pi}_{0,3}$ qui repère l'élé-%
ment courant $\gamma$ de $I\!\Gamma$. Ceci est lié à une représentation de la tour
de Teichmüller (en tant que groupoïde muni d'une opération de ``recol-%
lement'') par générateurs et relations, qui donnera en particulier
une présentation par générateurs et relations, au sens ordinaire, de
chacun des $\hat{T}_{g,\nu}$ (en tant que groupoïde profini). Ici, même pour
$g = 0$ (donc quand les groupes de Teichmüller correspondants sont
des groupes de tresses ``bien connus''), les générateurs et relations
connus à ce jour dont j'ai eu connaissance, me semblent inutilisables
tels quels, car ils ne présentent pas les caractères d'invariance et
de symétrie indispensables pour que l'action de $I\!\Gamma$ soit directement
lisible sur cette présentation. Ceci est lié notamment au fait que
les gens s'obstinent encore, en calculant avec des groupes fondamen-%
taux, à fixer un seul point base, plutôt que d'en choisir astucieu-%
sement tout un paquet qui soit invariant par les symétries de la si-%
tuation, lesquelles sont donc perdues en route. Dans certaines situa-%
tions (comme des théorèmes de descente à la Van Kampen pour groupes
fondamentaux) il est bien plus élégant, voir indispensable pour y com-%
prendre quelque chose, de travailler avec des groupoïdes fondamentaux
par rapport à un paquet de points bases convenable , et il en est cer-%
tainement ainsi pour la tour de Teichmüller. Il semblerait (incroyable,
mais vrai !) que la géométrie même du premier étage de la tour de
Teichmüller (correspondant donc aux ``modules'' soit pour des droites

\margin{7}
